\chapter{Sprint 3: Advanced Connectivity and API}
\chaptertoc

\section{Introduction}
Sprint 3 extended the system from a single local pipeline to a connected multi-source platform. The work introduced robust RTSP and ONVIF camera handling, a typed FastAPI backend-for-frontend (BFF), and runtime services for feed lifecycle and real-time event distribution.

\section{Sprint Backlog}
The Sprint 3 backlog addressed four integration tracks:
\begin{itemize}
    \item RTSP connectivity with authentication, transport control, reconnection, and snapshots.
    \item ONVIF discovery and stream resolution for IP camera onboarding.
    \item REST API endpoints for feed CRUD, metadata management, and diagnostics.
    \item Real-time delivery through websocket events and tiered video transport (WebRTC primary, MJPEG fallback, snapshots for zone workflows).
\end{itemize}

\section{Design}

\subsection{Component Diagram}
\begin{figure}[ht]
    \centering
    \fbox{\parbox[c][6cm][c]{0.85\textwidth}{\centering Placeholder for Sprint 3 component diagram}}
    \caption{Sprint 3 components for camera integration and backend services.}
    \label{fig:sprint3_component}
\end{figure}

The design separates responsibilities between camera adapters (RTSP/ONVIF), API contract models, runtime feed orchestration, and websocket broadcasting. This modularization allows independent evolution of source onboarding, analysis workers, and frontend-facing interfaces.

\subsection{Data Flow Diagram}
\begin{figure}[ht]
    \centering
    \fbox{\parbox[c][6cm][c]{0.85\textwidth}{\centering Placeholder for Sprint 3 data flow diagram}}
    \caption{Data flow from video feeds to API responses and WebSocket broadcasts.}
    \label{fig:sprint3_data_flow}
\end{figure}

The principal data paths are: (1) onboarding/test flows through REST endpoints, (2) feed-control flows through start/stop/restart operations, and (3) runtime event flows from worker JSONL streams to websocket broadcast and dashboard cache updates.

\subsection{Sequence Diagrams}
The implemented sequence for a running feed is:
frontend request $\rightarrow$ \texttt{/api/feeds/*} $\rightarrow$ \texttt{FeedRegistry} operation $\rightarrow$ worker process lifecycle change $\rightarrow$ websocket \texttt{feed\_status}/\texttt{metrics\_update} event.

\section{Implementation}
\subsection{RTSP Source Integration}
In \texttt{backend/src/rtsp\_camera.py}, the \texttt{RTSPCamera} class supports:
\begin{itemize}
    \item credential embedding with safe log redaction,
    \item TCP/UDP transport selection through OpenCV/FFMPEG options,
    \item connection testing and snapshot capture,
    \item automatic reconnection with bounded retries.
\end{itemize}

The source abstraction in \texttt{backend/src/video\_capture.py} routes local/webcam sources to \texttt{VideoStream} and RTSP URLs to \texttt{RTSPCamera}.

\subsection{ONVIF Discovery and Stream Resolution}
\texttt{backend/src/onvif\_client.py} implements WS-Discovery probes, ProbeMatch parsing, capability extraction, media-service resolution, and RTSP URI retrieval through SOAP requests with digest/basic authentication fallback. This enables dynamic IP camera onboarding without hardcoding stream URLs.

\subsection{FastAPI BFF and Typed Contract}
The API layer is implemented in \texttt{backend/src/api/app.py}, with typed models in \texttt{backend/src/api/models.py}. Major endpoint groups include:
\begin{itemize}
    \item feed management (create/list/start/stop/restart/delete),
    \item source onboarding (RTSP test/snapshot, ONVIF discovery/streams/test),
    \item metadata (establishments and cashiers),
    \item integration and health (webhook status/test, system health),
    \item live transport channels (websocket, transport capability endpoint, WebRTC signaling endpoint, and MJPEG stream endpoint).
\end{itemize}

\subsection{Transport Capability and WebRTC Signaling}
To reduce end-to-end viewing latency while preserving backward compatibility, Sprint 3 now exposes a capability-driven transport contract aligned with WebRTC signaling standards \cite{webrtc_w3c,ietf_rfc8830}.

The endpoint \texttt{/api/feeds/\{feed\_id\}/transport} returns per-feed \texttt{FeedTransportCapabilities} with both \texttt{webrtc} and \texttt{mjpeg} readiness fields. The frontend consumes this contract to decide transport routing:
\begin{itemize}
    \item \texttt{enabled}: feature-level availability,
    \item \texttt{ready}: runtime readiness for immediate use,
    \item \texttt{source\_mode}: \texttt{direct}, \texttt{annotated}, or \texttt{none},
    \item \texttt{reason}: machine-readable explanation when transport is unavailable.
\end{itemize}

When WebRTC is ready, the endpoint \texttt{/api/feeds/\{feed\_id\}/webrtc/offer} accepts a browser SDP offer and proxies it to MediaMTX WHEP, then returns the SDP answer \cite{mediamtx_docs}. The backend validates feed existence, running state, and RTSP-source eligibility before proxying.

Operational status semantics are explicit:
\begin{itemize}
    \item \texttt{200}: offer accepted and SDP answer returned,
    \item \texttt{404}: feed not found,
    \item \texttt{409}: feed not running, unsupported source, or transport not ready,
    \item \texttt{422}: invalid SDP payload,
    \item \texttt{502}/\texttt{503}: upstream or service-availability failures.
\end{itemize}

In \texttt{direct} mode, backend runtime can ensure MediaMTX path configuration through the control API before offer proxying. Environment-driven configuration (WHEP URL, control URL, timeout, and feature flag) keeps deployment adaptable without changing API contracts.

\subsection{Runtime Orchestration and Real-Time Events}
Runtime control resides in \texttt{backend/src/api/runtime.py}. \texttt{FeedRegistry} manages feed state, persistence sync, worker process lifecycle, and event dispatch. The worker command launches \texttt{src.main} in headless mode and writes event JSONL. \texttt{WebSocketHub} broadcasts standardized events:
\texttt{snapshot}, \texttt{feed\_status}, \texttt{metrics\_update}, \texttt{alert\_fired}, and \texttt{system\_warning}.

For video surfaces, \texttt{FeedFrameStreamManager} supports shared MJPEG buffering with capture-to-worker handoff continuity, preventing stream discontinuities when runtime mode changes.

\subsection{Persistence and Recovery}
\texttt{backend/src/database.py} provides SQLite persistence for feed configuration, metadata, video sessions, and archived alerts. On API startup, persisted feeds are restored and stale running sessions are safely recovered to a stopped state.

\section{Conclusion}
Sprint 3 delivered the system connectivity layer and formal service contract required by the web application. With source onboarding, runtime orchestration, typed API/websocket channels, and capability-driven WebRTC signaling in place, the project was ready for full web UI integration and operational workflows.
